\documentclass[12pt]{article}

% ── Packages ──────────────────────────────────────────────────────────────
\usepackage[margin=1in]{geometry}
\usepackage{amsmath, amssymb}
\usepackage{graphicx}
\usepackage{booktabs}
\usepackage{natbib}
\usepackage{hyperref}
\usepackage{xcolor}
\usepackage{soul}          % \hl{} for TODO highlights
\usepackage{enumitem}
\usepackage{array}
\usepackage{caption}
\usepackage{subcaption}
\usepackage{multirow}
\usepackage{lineno}        % line numbers for review
\linenumbers

% ── TODO command ──────────────────────────────────────────────────────────
\newcommand{\todo}[1]{\hl{\textbf{[TODO: #1]}}}
\newcommand{\note}[1]{{\color{blue}\textit{[Note: #1]}}}

% ── Title block ───────────────────────────────────────────────────────────
\title{\textbf{Latent Quality Risk in Generic Drug Manufacturing:\\
  LLM-Extracted FDA 483 Signals and Drug Shortage Prediction}}

\author{
  Amirreza Sahebi\thanks{Corresponding author.
    Poole College of Management, North Carolina State University.
    Email: \texttt{asahebi@ncsu.edu}} \\
  \and
  \todo{Co-authors}
}

\date{\today \\ \vspace{4pt}
  \textit{Preliminary draft — prepared for INFORMS Annual Meeting 2026.
  Please do not cite or circulate without permission.}}

% ══════════════════════════════════════════════════════════════════════════
\begin{document}
\maketitle

% ── Abstract ──────────────────────────────────────────────────────────────
\begin{abstract}
Generic drug shortages pose significant risks to patient safety and public
health, yet their prediction remains challenging due to the latent nature of
manufacturing quality failures.
FDA Form 483 inspectional observations—letters issued to drug manufacturing
facilities documenting specific violations—contain rich qualitative information
about the nature, severity, and root causes of quality deficiencies, but this
information has not been systematically quantified for risk prediction.
We apply large language model (LLM) extraction to \todo{N} Form 483
observations from \todo{N} facilities manufacturing 14 generic APIs tested
by Valisure, converting free-text observations into structured signals:
violation severity, root-cause type, remediation quality, and seven binary
risk flags (repeat, systemic, patient risk, data integrity, contamination,
documentation, investigation failure).
We aggregate these signals to facility-level Text Risk Indices (TRI, SCRI,
IRWI, QCI) and test whether they improve prediction of Valisure-measured
quality failures and University of Utah drug shortage onset.
Our results suggest that \todo{key finding}.
The approach is fully reproducible from publicly available FDA data and offers
a scalable framework for real-time manufacturing quality surveillance.
\end{abstract}

\textbf{Keywords:} drug shortage prediction, FDA Form 483, manufacturing quality,
large language model, text analysis, pharmaceutical supply chain

\newpage

% ══════════════════════════════════════════════════════════════════════════
\section{Introduction}
\label{sec:intro}

Generic drugs account for roughly 90\% of U.S. prescriptions filled
\citep{fda2023generic}, yet generic drug shortages are an enduring public
health problem that disrupts patient care, inflates costs, and strains
hospital operations \citep{fox2014shortage}.
Shortages arise from an interplay of demand concentration, supply chain
fragility, and latent manufacturing quality failures that may not surface
until facilities fail FDA inspections or recall products.

The FDA's regulatory data—inspection outcomes (OAI/VAI/NAI classifications),
warning letters, recalls, and import refusals—provide structured signals of
quality risk, but they are lagging indicators: they reflect quality failures
that have already materialized.
FDA Form 483 letters are issued at the \emph{conclusion} of an inspection and
document specific deficiencies before they escalate to warning letters or
recalls.
Form 483 observations thus constitute a leading signal of manufacturing
quality degradation, but their unstructured narrative format has prevented
systematic quantitative use.

\todo{2--3 sentences on prior literature: structured regulatory signals for
shortage prediction (Fox et al., Pauwels et al., etc.); FDA text analysis
work (if any); LLM for regulatory/clinical text.}

We make three contributions:
\begin{enumerate}[noitemsep]
  \item We develop a validated LLM extraction pipeline that converts Form 483
        observation text into structured, reproducible risk signals at the
        observation level.
  \item We construct four facility-level Text Risk Indices—TRI, SCRI, IRWI,
        and QCI—aggregating these signals and empirically characterize their
        distributions across \todo{N} generic drug manufacturing facilities.
  \item We test whether these indices improve prediction of independently
        verified quality failures (Valisure dissolution testing) and drug
        shortage onset (University of Utah shortage database) relative to
        structured regulatory features alone.
\end{enumerate}

% ══════════════════════════════════════════════════════════════════════════
\section{Background and Data}
\label{sec:data}

\subsection{FDA Form 483 Observations}

FDA Form 483 is issued to a facility's management at the close of an
inspection when an investigator observes conditions that may violate the
Food, Drug, and Cosmetic Act or its implementing regulations (21 CFR Parts
210–211).
Each letter contains one or more numbered observations, each describing a
specific deficiency.
Observations are written by FDA investigators in narrative form; they do not
follow a templated structure beyond the opening CFR-paraphrase sentence.

We obtained \todo{N} Form 483 PDFs from FDA's public disclosure portal,
covering \todo{N} facilities and the period \todo{year range}.
PDFs were converted to text via OCR (Tesseract 200 DPI) and parsed into
individual observations using regex splitting on ``OBSERVATION N'' headers,
yielding \textbf{622 observations} across \textbf{38 facilities} after
quality filtering.

\subsection{Valisure Quality Outcomes}

Valisure is an independent pharmaceutical testing laboratory that publicly
reports dissolution testing, nitrosamine contamination (NDMA, DMF), and
quality failure results for commercially available generic drugs.
We use Valisure test results for 14 APIs tested in 2020, 2022, and 2024 as
our primary quality outcome variable.
These results serve as independently verified ground truth for manufacturing
quality, free of the endogeneity concerns of FDA inspection outcomes.

\subsection{Drug Shortage Data}

Drug shortage onset and duration data are drawn from the University of Utah
Drug Information Service (UUtah) shortage database, which tracks generic drug
shortages from \todo{year range}.
We define shortage onset as a binary indicator for a given API--year pair.

\subsection{Sample}

\todo{Describe the 14 APIs, FEI-to-drug mapping via NDC-FEI crosswalk,
panel structure (drug-year or FEI-year), and any sample selection rules.}

% ══════════════════════════════════════════════════════════════════════════
\section{Methods}
\label{sec:methods}

\subsection{LLM Extraction Pipeline}

For each Form 483 observation we send the cleaned observation text to
\texttt{gpt-4o-mini} (OpenAI) with a structured JSON schema prompt.
The prompt instructs the model to return exactly 15 fields:

\begin{itemize}[noitemsep]
  \item \textbf{Categorical}: violation category (8 values), severity tier
        (High/Moderate/Low), root-cause type (Capital/Cultural/Mixed/Unclear),
        remediation signal (Strong/Partial/Weak/None).
  \item \textbf{Binary flags} (7): repeat finding, systemic scope, patient
        risk, data integrity failure, contamination/sterility risk,
        documentation deficiency, investigation failure.
  \item \textbf{Qualitative}: severity rationale, root-cause rationale,
        verbatim evidence quote, confidence score (0--1).
\end{itemize}

The evidence quote is subject to an \emph{evidence guard}: the returned quote
must appear verbatim in the source text; otherwise it is cleared.
Extraction is idempotent: already-scored observations are skipped on re-run.

\paragraph{Validation.}
We manually reviewed a random sample of \todo{N} observations and compared
LLM classifications to independent human raters.
Inter-rater agreement: \todo{Cohen's $\kappa$ or \%}.

\subsection{Facility-Level Aggregation}

We aggregate observation-level signals to the facility (FEI) level using
the following composite indices, each bounded $[0, 100]$:

\begin{align}
\text{TRI}_i &= \left[
    0.35 \cdot s^H_i
  + 0.20 \cdot s^M_i
  + 0.20 \cdot (1 - r^S_i)
  + 0.15 \cdot f^{\text{rep}}_i
  + 0.10 \cdot f^{\text{sys}}_i
\right] \times 100 \label{eq:tri} \\
\text{SCRI}_i &= \left[
    0.50 \cdot f^{\text{contam}}_i
  + 0.30 \cdot f^{\text{contam,rx}}_i
  + 0.20 \cdot s^H_i
\right] \times 100 \label{eq:scri} \\
\text{IRWI}_i &= \left[
    0.40 \cdot r^N_i
  + 0.35 \cdot f^{\text{invest}}_i
  + 0.25 \cdot r^W_i
\right] \times 100 \label{eq:irwi} \\
\text{QCI}_i &= \left[
    0.40 \cdot f^{\text{sys}}_i
  + 0.35 \cdot f^{\text{rep}}_i
  + 0.25 \cdot c^{\text{cult}}_i
\right] \times 100 \label{eq:qci}
\end{align}

where $s^H_i$, $s^M_i$ are severity High/Moderate shares; $r^S_i$, $r^N_i$,
$r^W_i$ are remediation Strong/None/Weak shares; $f^{\cdot}_i$ are LLM flag
shares; $f^{\text{contam,rx}}_i$ is the regex contamination share; and
$c^{\text{cult}}_i$ is the cultural root-cause share, all for facility $i$.

\subsection{Manufacturing Quality Risk Index (MQRI)}

The MQRI combines regulatory enforcement features and market quality signals:

\begin{equation}
  \text{MQRI}_i = \left( 0.70 \cdot D^{\text{reg}}_i
                         + 0.30 \cdot D^{\text{mkt}}_i \right) \times 100
\end{equation}

where $D^{\text{reg}}$ includes cumulative OAI count, OAI-predictive CFR
inspections, posted 483 citations, warning letters, inspection gap, import
refusals, and the contamination NLP flag; $D^{\text{mkt}}$ includes Class I
and all-class recall counts.
Feature weights are derived from Spearman $|\rho|$ with OAI rate across
$n = 127$ facilities.
In this paper we extend the MQRI by adding TRI as a regulatory domain feature
with weight derived from \todo{correlation with OAI rate or Valisure}.

\subsection{Shortage Prediction Model}

\todo{Describe m09 model: logistic regression / XGBoost / LASSO at drug--year
level, feature set from m07 panel, train/test split, evaluation metric (AUC).
Then describe the extension: adding TRI/SCRI via FEI→drug bridge.}

% ══════════════════════════════════════════════════════════════════════════
\section{Results}
\label{sec:results}

\subsection{Descriptive Characterization of Form 483 Observations}

Table~\ref{tab:obs_desc} summarizes the LLM extraction results across
622 observations from 38 facilities.

\begin{table}[h!]
\centering
\caption{Observation-level signal distribution ($n = 622$ observations,
         38 facilities)}
\label{tab:obs_desc}
\begin{tabular}{lrr}
\toprule
\textbf{Signal} & \textbf{Count} & \textbf{Share (\%)} \\
\midrule
\multicolumn{3}{l}{\textit{Severity tier}} \\
\quad High      & 435 & 69.9 \\
\quad Moderate  & 174 & 28.0 \\
\quad Low       &  13 &  2.1 \\
\midrule
\multicolumn{3}{l}{\textit{Primary violation domain}} \\
\quad Production Controls  & 170 & 27.3 \\
\quad Lab Controls         & 162 & 26.1 \\
\quad Quality System       & 123 & 19.8 \\
\quad Buildings/Equipment  &  71 & 11.4 \\
\quad Records/Reports      &  61 &  9.8 \\
\quad Personnel            &  22 &  3.5 \\
\quad Other                &  13 &  2.1 \\
\midrule
\multicolumn{3}{l}{\textit{Root-cause type}} \\
\quad Cultural  & 279 & 44.9 \\
\quad Mixed     & 242 & 38.9 \\
\quad Capital   &  92 & 14.8 \\
\quad Unclear   &   9 &  1.4 \\
\midrule
\multicolumn{3}{l}{\textit{Remediation signal}} \\
\quad None    & 418 & 67.2 \\
\quad Weak    & 112 & 18.0 \\
\quad Partial &  91 & 14.6 \\
\quad Strong  &   1 &  0.2 \\
\bottomrule
\end{tabular}
\end{table}

\paragraph{Key finding 1: Severity is predominantly high.}
Nearly 70\% of 483 observations in our sample are classified as High severity
(direct patient-harm risk or gross deviation from CGMP).
This reflects the facility composition of our sample—facilities that
manufactured drugs subsequently tested by Valisure, which may be a
non-random subset with higher historical inspection activity.

\paragraph{Key finding 2: Cultural root causes dominate.}
83.8\% of observations are classified as Cultural or Mixed root cause,
suggesting that the predominant drivers of CGMP violations in this sample
are management, training, and quality-culture failures rather than
capital/equipment gaps.
This has direct implications for corrective action: capital investment alone
is insufficient.

\paragraph{Key finding 3: Remediation is rarely mentioned.}
67.2\% of observations mention no corrective action whatsoever, and only
one observation (0.2\%) references a strong corrective action.
\note{This is partly structural: 483 observations are investigator-written
findings; firm responses are submitted separately.
However, Partial/Weak mentions do appear when investigators note commitments
made on-site. The absence is itself informative.}

\subsection{Semantic Lift: LLM vs.\ Regex}

Table~\ref{tab:lift} compares LLM flag rates to deterministic regex flags
on the same observations.

\begin{table}[h!]
\centering
\caption{Semantic lift: LLM vs.\ regex flag rates (mean across 38 FEIs)}
\label{tab:lift}
\begin{tabular}{lrrr}
\toprule
\textbf{Flag} & \textbf{Regex share} & \textbf{LLM share}
              & \textbf{LLM-only lift} \\
\midrule
Patient risk   & 19.4\% & $\sim$90\% & \textbf{71.1\%} \\
Systemic scope & 26.3\% & $\sim$95\% & \textbf{68.9\%} \\
Contamination  & 30.2\% & $\sim$57\% & 27.1\% \\
Data integrity &  6.1\% & $\sim$19\% & 13.1\% \\
Investigation  & 39.9\% & $\sim$43\% &  3.2\% \\
Documentation  & 83.6\% & $\sim$87\% &  4.7\% \\
Repeat finding &  4.3\% & $\sim$9\%  &  5.0\% \\
\bottomrule
\end{tabular}
\end{table}

The patient risk and systemic flags show the largest semantic lift:
the LLM identifies \emph{patient risk implications} and \emph{facility-wide
scope} in 71\% and 69\% of observations, respectively, where keyword
rules fail.
This is because these inferences require contextual reasoning
(e.g., sterile injectable $\Rightarrow$ patient risk; ``your firm's''
$\Rightarrow$ systemic) that regex cannot replicate.

\subsection{Facility-Level Text Risk Indices}

\todo{Table of top/bottom 10 FEIs by TRI with firm name (or anonymized).
Add: correlation of TRI with severity\_score from node\_summary,
correlation of TRI with n\_oai, correlation with MQRI.}

\subsection{Predictive Validity: TRI vs.\ Valisure Outcomes}

\todo{Run correlation: TRI vs. Valisure dissolution difference factor (2024).
Run correlation: SCRI vs. Valisure NDMA/DMF.
Compare to MQRI-only baseline.}

\subsection{Predictive Validity: TRI in Shortage Prediction}

\todo{Run m09 with and without TRI/SCRI features.
Report AUC with 95\% CI, precision-recall curve.
Feature importance.}

% ══════════════════════════════════════════════════════════════════════════
\section{Discussion}
\label{sec:discussion}

\paragraph{Form 483 as a leading quality signal.}
\todo{Discussion of why 483 text adds information beyond structured signals.}

\paragraph{Root cause composition and policy implications.}
The predominance of Cultural and Mixed root causes suggests that FDA
enforcement actions alone are insufficient remedies.
Facilities that receive OAI classifications repeatedly are more likely
characterized by management culture failures than capital deficiencies,
implying that CAPA programs targeting equipment or procedures may be
systematically underspecified.

\paragraph{Limitations.}
\begin{itemize}[noitemsep]
  \item \textit{Coverage}: only 38 of 129 facilities in our FEI sample have
        associated 483 PDFs, limiting the reach of text-derived features.
  \item \textit{LLM reliability}: while precision on manually reviewed
        observations is \todo{\%}, the model occasionally over-triggers
        contamination flags for training observations.
  \item \textit{Selection}: Valisure tests a non-random subset of generics;
        our results may not generalize to lower-scrutiny drug classes.
  \item \textit{Causality}: we estimate predictive associations, not causal effects.
\end{itemize}

% ══════════════════════════════════════════════════════════════════════════
\section{Conclusion}
\label{sec:conclusion}

We demonstrate that LLM-extracted signals from FDA Form 483 observations
provide structured, reproducible, and predictively informative features for
manufacturing quality risk assessment.
The proposed Text Risk Index and related indices capture dimensions of quality
risk—severity composition, root-cause type, and remediation quality—that are
absent from existing structured regulatory databases.
Integrating these features into the MQRI and shortage prediction models
\todo{improves/does not improve} predictive performance.
The pipeline is fully automated and can be applied prospectively to the
FDA's public 483 disclosure portal, enabling near-real-time quality
surveillance for drug shortage early warning.

% ══════════════════════════════════════════════════════════════════════════
\bibliographystyle{agsm}
\bibliography{references}

% ══════════════════════════════════════════════════════════════════════════
\appendix

\section{LLM Flag Definitions}
\label{app:flags}

\paragraph{Extraction model.} \texttt{gpt-4o-mini} with structured JSON output
schema (strict mode). Temperature = 0 (deterministic). Max tokens = 2000.

\paragraph{Evidence guard.} The \texttt{evidence\_quote} field must appear
verbatim in the source observation text. If the model paraphrases, the field
is cleared to empty string rather than accepting an inaccurate citation.

\paragraph{Flag definitions} are as specified in the extraction prompt
(see \texttt{01\_extract\_observation\_signals.py}):

\begin{description}[leftmargin=2em]
  \item[\texttt{repeat\_flag\_llm}] True only when the observation explicitly
    states this is a repeat or previously-cited finding from a prior inspection.
    Not true for repeated examples within the same current observation.
  \item[\texttt{systemic\_flag\_llm}] True when the observation describes
    facility-wide, multi-process, multi-product, or quality-system-level failure.
  \item[\texttt{patient\_risk\_flag\_llm}] True when the violation could
    directly affect patient safety. Does not require confirmed patient harm.
  \item[\texttt{data\_integrity\_flag\_llm}] True only for explicit data
    trustworthiness failures (falsification, backdating, audit-trail problems).
    Not true for ordinary incomplete documentation.
  \item[\texttt{contamination\_flag\_llm}] True for actual contamination or
    clear contamination-control risk including sterility assurance failures.
  \item[\texttt{documentation\_flag\_llm}] True when inadequate documentation
    is a \emph{central} finding.
  \item[\texttt{investigation\_flag\_llm}] True only for an explicit failed,
    missing, or inadequate investigation of a concrete event.
\end{description}

\section{Composite Index Weights}
\label{app:weights}

Weights for TRI, SCRI, IRWI, and QCI are specified in
\texttt{02\_aggregate\_fei\_features.py} and reproduced in
Equations~\eqref{eq:tri}--\eqref{eq:qci}.
\todo{Add empirical derivation table once correlations with OAI rate are run.}

\end{document}
